\documentclass[11pt]{article}
\usepackage[margin=1in]{geometry}
\usepackage{amsmath,amssymb,amsthm}
\usepackage{enumitem}

\newcommand{\Om}{\Omega}
\newcommand{\dT}{\delta_T}
\newcommand{\A}{\mathcal A}
\newcommand{\Sph}{\mathbb S}
\newcommand{\R}{\mathbb R}
\newcommand{\Ht}{\widehat t}
\newcommand{\Et}{E_{\widehat t}}

\title{Plan 3 Graph-Reduction Recheck}
\author{}
\date{May 2026}

\begin{document}
\maketitle

\section{Purpose}

I rechecked the Plan 3 material for the specific claim that the reduction from
a good torsion level to a radial graph is already closed.  The conclusion is
negative in the literal sense: the repo contains several real ingredients for
such a reduction, but no complete proof that the stopped level is, or can be
repaired at \(O(\dT)\) cost into, a small \(L^\infty\) radial graph.

The closest-looking closures are:
\begin{enumerate}[label=\arabic*.]
\item \texttt{C-assembled-proof-steps/proof-step2.md}, especially Sections
7--9, plus \texttt{B-routeB-weinberger/verify-chainB.md};
\item \texttt{C-assembled-proof-steps/proof-step4.md}, which explains how an
\(L^\infty\) Serrin output would enter the graph theorem;
\item \texttt{A-strategyA-graph/expA-anygraph-monotonicity.md} and the
Strategy B notes;
\item \texttt{B-routeB-weinberger/tentacle-model.md};
\item the Plan 2 metric/interval route, which replaces graph coordinates by
an \(L^1\)-modulo-translations metric trace.
\end{enumerate}

None of these proves the desired graph-entry or radialization theorem as
written.

\section{The strongest apparent closure: Chain B}

The text in \texttt{proof-step2.md} says that the graph substatement is
unnecessary.  It proposes the chain
\[
   \dT(D)
      \gtrsim |D|^{-1}
      \int_D \left|D^2u_D+\frac1nI\right|^2
      \quad\Longrightarrow\quad
   \A(D)^2
      \lesssim
      \int_D \left|D^2u_D+\frac1nI\right|^2,
\]
for \(D=\Et\).  This would indeed bypass graph entry.  The problem is that
the first inequality is false, even on smooth volume-normalized near-balls.

Let \(Y_k\) be an \(L^2(\Sph^{n-1})\)-normalized spherical harmonic of degree
\(k\ge 2\) and
\[
   D_{\varepsilon,k}
      =\{r<1+\varepsilon Y_k+O_k(\varepsilon^2)\},
\]
with the \(O_k(\varepsilon^2)\) constant mode enforcing volume.  The elementary
ball linearization gives
\[
   \dT(D_{\varepsilon,k})
      =
      \frac{\varepsilon^2}{2n^2}(k-1)+O_k(\varepsilon^3),
\]
whereas
\[
   \int_{D_{\varepsilon,k}}
      \left|D^2u_{D_{\varepsilon,k}}+\frac1nI\right|^2
      =
      \frac{\varepsilon^2}{n^2}(k-1)k(2k+n-2)+O_k(\varepsilon^3).
\]
Thus
\[
   \frac{\dT(D_{\varepsilon,k})}
        {\int_{D_{\varepsilon,k}}|D^2u_{D_{\varepsilon,k}}+I/n|^2}
      \to
      \frac{1}{2k(2k+n-2)}
      \to 0.
\]
This is the calculation recorded in
\texttt{PLAN3\_N2\_SINGLE\_LEVEL\_RECHECK.tex}.  It kills the unweighted
Hessian coercivity used by Chain B.  Since the counterexamples are smooth
small radial graphs, the failure is not caused by tentacles, disconnectedness,
or rough boundary.

\section{Why the auxiliary \(H^2\)-to-geometry text does not repair this}

Even if a small unweighted \(H^2\) defect were somehow available, the
conversion argument in \texttt{proof-step2.md} Section 7 is not a proved
graph-reduction theorem.

The proof asserts that after removing a bad set \(Z\), the remaining domain is
a uniformly John/Lipschitz union foliated by nondegenerate levels, and that
local harmonic estimates can be chained to one constant vector.  The inputs
cited there are:
\[
   \mathcal H^{n-1}\{|\nabla u|\le s\}\lesssim D_H\,s
\]
and a model tube estimate saying that the unweighted Hessian defect sees thin
tubes.  These inputs do not imply a uniform John constant, a single connected
bulk component, or a quantitative no-fjord/no-hole statement.  A small-volume
exceptional set can still disconnect the bulk or leave same-ray multiple
crossings.

The final step also asserts that \(H^1(D)\)-closeness of \(u_D\) to a radial
quadratic implies \(L^1\)-closeness of \(D\) to a ball by a coarea/trace
argument.  That implication needs quantitative nondegeneracy in a full collar
around the zero level, not merely small surface measure of a low-gradient set
on the level.  In particular, it does not give a star-shaped graph.

Therefore the parenthetical line in \texttt{proof-step2.md} saying that
``\(H^2\) defect plus volume exclusion gives a genuine star-shaped graph with
small \(H^{1/2}\) norm'' is an assertion, not a proof.  It is also not
load-bearing once Chain B is removed.

\section{The Step 4 graph passage assumes the missing input}

\texttt{proof-step4.md} correctly says that if Step 2 produced a smooth
closed level with pointwise Serrin control, nondegeneracy, and genuine
star-shaped graph structure, then the implicit-function/overdetermined map
would put it in a perturbative graph framework.  But this is conditional on
the missing Step-2 output:
\[
   \bigl\||\nabla u|-c\bigr\|_{L^\infty(\Sigma_{\Ht})}\ll1
   \quad\hbox{and}\quad
   \Sigma_{\Ht}\hbox{ is a small graph.}
\]
The corrected Step 1 gives a good level with an \(O(\dT)\) removed layer and
integrated \(D_H+D_I\) control.  It does not give this pointwise Serrin bound,
a uniform gradient floor, a \(C^1\) graph, or one radial crossing.  Thus Step 4
is a valid downstream block, not a graph-entry proof.

\section{Any-graph/Strategy B does not supply the reduction}

\texttt{expA-anygraph-monotonicity.md} contains a useful warning: the naive
radial homotopy monotonicity proof fails because higher-order shape terms have
no sign on finite non-small graphs.  It also points to a Pohozaev/\(P\)-function
route for star-shaped graph domains.

That route cannot be used as written for the present reduction.  First, it
assumes star-shaped graph geometry, which is exactly what we are trying to
derive.  Second, its advertised estimate
\[
   \dT(D)\gtrsim \int_D |D^2u_D+I/n|^2
\]
is the same unweighted Hessian coercivity contradicted by the high-frequency
near-ball calculation above.  Taking \(\varepsilon\ll k^{-1}\) keeps
\(D_{\varepsilon,k}\) star-shaped with small \(C^1\) norm, while the ratio of
\(\dT\) to the Hessian defect still tends to zero as \(k\to\infty\).

\section{What the \(n=2\) tentacle computation really gives}

\texttt{tentacle-model.md} is genuinely useful.  It shows that the specific
``move inward to escape an outward thin tentacle'' bookkeeping is favorable in
dimension two.  For a tube of width \(w\) and length \(L\),
\[
   \dT \simeq wL,
   \qquad
   \hbox{escape layer}\simeq w^2,
\]
so in \(n=2\) the layer is paid for when \(L\gg w\).  This removes one serious
model obstruction.

But it is not a graph-entry theorem.  It does not prove an arbitrary-geometry
decomposition of every outward non-graph component of an analytic torsion
level into chargeable fingers.  It also says nothing by itself about inward
fjords, holes, or same-ray multiple crossings.

\section{What Plan 2 contributes}

The Plan 2 radial slicing and metric-trace machinery is the strongest
replacement for graph geometry in the repo.  It gives:
\begin{enumerate}[label=\arabic*.]
\item finite-perimeter level identities with exact \(D_H+D_I\);
\item same-center isoperimetric/radial trace control in \(L^1\)-type metrics;
\item bad-ray multiplicity absorption once an annular bracket is known;
\item a plausible interval route at the correct \(H^{1/2}\) scale.
\end{enumerate}

This is not the same as graph entry.  It allows a small set of bad rays and
works in \(L^1\)/metric distance, whereas the Plan 3 graph theorem needs a
single-valued \(L^\infty\) radial graph or an \(O(\dT)\)-cost repair to one.

Moreover the current Plan 2 interval route still records an explicit remaining
conditional input.  In \texttt{WIP\_WeightedMetricTrace.tex} the bad-level
kinetic hypothesis asks for
\[
   \int_{B_I}|\dot F_\rho|_{\mathcal X}^2\,d\mu\le C\dT,
   \qquad
   \int_{B_{\tau_0}}|\dot F_\rho|_{\mathcal X}^2\,d\rho\le C\dT.
\]
The patch note \texttt{PLAN2\_INTERVAL\_ROUTE\_BAD\_KINETIC\_PATCH.md}
closes the bad-\(I\) part and the slow-profile part on \(G_I\), but leaves
\[
   \int_{B_{\tau_0}\cap B_I}|\dot F_\rho|_{\mathcal X}^2\,d\rho
      \le C\dT
\]
open.  Equivalently, with the crude deformation bound, it would suffice to
prove
\[
   \int_{B_{\tau_0}\cap B_I}D_I(t(\rho))\,d\rho\le C\dT.
\]
The existing identities control only the weighted integral
\(\int D_I(t(\rho))w(\rho)\,d\rho\), and on \(B_{\tau_0}\) the weight
\(w=-t'(\rho)\) is small.  Thus the interval route is promising, but it does
not close the Plan 3 graph reduction.

\section{What survives for the hybrid Plan 3 route}

The following pieces are sound and should be retained.

\begin{enumerate}[label=\arabic*.]
\item The corrected Step 1 good-level extraction: one can stop at a regular
near-boundary level with \(O(\dT)\) removed volume and integrated
\(D_H+D_I\) control.
\item The superlevel torsion transfer and the final symmetric-difference
transfer back to \(\Om\).
\item The low-gradient surface estimate
\[
   \mathcal H^{n-1}(\{|\nabla u|\le s\}\cap\Sigma_{\Ht})
      \lesssim D_H(\Ht)\,s.
\]
It is a real quantitative fact, but it is only measure control.
\item The direct radial-graph theorem in
\texttt{H12\_ASYMMETRY\_DIRECT\_PROOF.tex}: if a domain is a genuine
small-amplitude radial graph, no slope or curvature bound is needed to get
\(\A^2\lesssim \dT\).
\item The \(n=2\) outward-tentacle model: outward fingers should be chargeable
at the correct \(O(\dT)\) scale in dimension two.
\item The capacity identity in
\texttt{PLANAR\_RADIALIZATION\_PROOF\_ATTEMPT.tex}: if an inward missing set
lies at positive torsion height in the filled comparison domain, then filling
it has a torsion cost at least \(a^2\operatorname{cap}(K)\).
\end{enumerate}

These pieces strongly suggest the right missing theorem, but they do not prove
it.

\section{The exact remaining graph-reduction theorem}

The honest target is the \(O(\dT)\) radialization lemma:

\medskip
\noindent
\textbf{Planar radialization lemma.}
Let \(n=2\), let \(E=\Et\) be the corrected Step-1 good level, and assume
\(\dT(\Om)\) is small in the bounded normalized class.  Then, after translation
and rescaling, there is a set
\[
   G=\{re^{i\theta}:0<r<1+\varphi(\theta)\},
   \qquad
   \|\varphi\|_{L^\infty}\le \eta_0,
\]
such that
\[
   |E\Delta G|\le C\dT(\Om),
   \qquad
   \dT(G)\le C\dT(\Om).
\]

If this lemma is proved, the existing radial-graph theorem gives
\(\A(G)^2\le C\dT(G)\), the \(O(\dT)\) repair gives
\(\A(E)^2\le C\dT(\Om)\), and Step 5 transfers back to \(\Om\).

The missing proof should split into three pieces:
\begin{enumerate}[label=\arabic*.]
\item an annular/core theorem putting the good level in a fixed thin annulus
and giving positive torsion height on the pieces to be filled;
\item a planar capacity/modulus fill lemma for inward fjords, holes, and
same-ray missing intervals;
\item an arbitrary-geometry outward pruning and bad-ray radialization lemma,
using the \(n=2\) \(D_H\)-tentacle scaling plus Plan 2 bad-ray multiplicity.
\end{enumerate}

\section{Verdict}

I think the memory is pointing to real material, but it closes less than the
literal graph-reduction step.  The repo already contains:
\[
   \hbox{good level}
   + \hbox{metric/radial \(L^1\) controls}
   + \hbox{\(n=2\) outward-tentacle bookkeeping}
   + \hbox{capacity mechanism for fillable inward defects}.
\]
What it does not contain is the theorem that turns these into either a genuine
small \(L^\infty\) radial graph or an \(O(\dT)\)-cost radialized replacement.
The apparent closures in Plan 3 use a false Hessian inequality, assume the
graph entry, or prove only model/metric statements.

\end{document}
